\section{Cameras and Visual Detectors}
\label{sec:cameras}

\subsection{Cameras}
The system takes input from two cameras -- one Logitech camera (the top camera) elevated in a fixed position above the robot and one RealSense depth camera attached to the robot right above the gripper. The top camera is primarily used for identifying blocks in the "pile" -- a section of randomly assorted blocks to the right of the robot that it can use to build the structure --  to be used to build the structure as well as identifying whether the robot accidentally picked up multiple blocks at once. The RealSense camera is used primarily when building the structure to determine the structure state as it can determine the z coordinate of the blocks in the structure, and can therefore tell us which layer the blocks in the structure are in.

\subsection{Calibration}
\subsubsection{Aruco Markers}
We produce a mapping of pixel-to-world from the top camera frames with the positioning of Aruco markers at the corners of the workspace both at the start of the program and periodically (typically every 100 seconds) as the robot runs in case the table shifts. The detector knows the world coordinates of the Aruco markers, so it uses the detected pixel coordinates of those markers to create a affine transform matrix to convert between pixel and world coordinates.

\subsubsection{XY Calibration}
\label{sec:xy_calib}
We calibrate the robot's kinematics errors in the $xy$ plane before the start of construction by having the robot move to various positions in the workspace, detecting the position of a blue marker on the robot's camera mount at those points with the top camera, and storing the offset between the expected position (i.e. the task command sent to the robot) and the actual position of the tip. These offsets are then used by the inverse kinematics solver to offset the task coordinates by an amount computed by interpolating the offsets (Figure~\ref{fig:xy_error}) between the different points that the calibration was done at.

\noindent It's important to note that since the blue mark on the camera mount is not flat with the table, we implemented a ray-traced-based way to convert pixel coordinates to world coordinates with accurate $z$ mapping.

\subsubsection{Z Calibration}
\label{sec:z_calib}
We calibrate the robot's kinematics errors on the $z$-axis (caused primarily by imperfect gravity compensation) before the start of construction by having the robot move along the radial direction at different $z$ positions and detecting the height of the table at those positions with the end effector camera. As the tip moves farther out in the radial direction, the table appears higher in world coordinates even though it is meant to be at a constant height of 0. This is because the imperfect gravity compensation has a greater effect as the robot extends out further (Figure~\ref{fig:z_error}). We store these table heights at different radii and apply an offset to task coordinates before initiating the inverse kinematics solve, similar to the $xy$ calibration.

\subsubsection{Placement Grid Calibration}
The placement grid calibration determines the world-frame positions of the structure's base grid corners using the overhead camera. While the grid corner coordinates are initially measured by hand, this step re-detects them through the same camera and calibration pipeline used for the XY kinematics calibration. This ensures that the camera's understanding of where the grid is uses the same pixel-to-world transform as the kinematics error correction, preventing errors from different subsystems (camera calibration vs.\ kinematics) from compounding.

\subsection{Block Detection: Pile}
To detect potential blocks from the pile we use a simple contour detection algorithm for the top camera as the blocks we are detecting are always squares. We first filter by color, using HSV limits to define the colors corresponding to the blocks that we use. We then use the OpenCV library to isolate the contours of objects within those HSV limits. We then use a custom edge-aligned bounding box method, which identifies the edges of the contour and the angle of the longest edge, rotates the contour by this angle, computes the axis aligned bounding box of this rotated contour, rotates this axis aligned bounding box back to the countour's angle, and returns it. This method finds a bounding box for the contour in a manner that aligns with the actual blocks within it, and therefore allows us to identify where the blocks are within the contour. Here we have three cases: either the bounding box is smaller than the size of a typical block, in which case we discard it; the bounding box is approximately the size of a block, in which case we just detect it as a single block, or it is the size of at least two or more typical blocks, in which case we split it up into its internal blocks, which we mark as connected.

\noindent To split it up we use a custom "rectangle splitting algorithm", which when given one of these large bounding rectangles (detected by either the top or end effector camera), will split this large bounding rectangle into individual detected blocks and determined which blocks within that rectangle are connected. The algorithm works by taking the bounding rectangle and splitting it into a grid based on the expected dimensions of a block in pixel space at the height of the camera. It then checks whether each square in that grid is occupied by a block and designates the corresponding block as connected to the other detected blocks adjacent to it in the grid if it is occupied. This allows us to detect which faces of any detected block are connected.

\subsection{Block Detection: Structure}
To determine the current completeness state of the structure, we use the point cloud (points specified in xyz with RGB color) generated by end-effector depth camera to accumulate an occupancy grid (which is explained later) which we can query to detect blocks in 3D. For each layer (z modulo block size), we query the occupancy map and filter by color using HSV limits. This gives us a layer-wise + color-wise binary map which we can use the exact algorithm as for the block detection from pile (equivalent to layer=0) to detect blocks from. Importantly, the structure detector operates in grid-relative coordinates: detected block positions are mapped relative to the base grid corners (via an affine transform matrix) rather than in absolute world coordinates. This decouples kinematics errors from the world-to-end-effector camera transform errors, as both the grid registration and the block detections pass through the same camera pipeline.

\subsection{Grip \& Connection Detection}
In addition to object detection, we are able to validate whether we grabbed blocks properly with the two cameras. Using the end effector camera, we detect the depth of the grasped block between the gripper. This allows us to tell whether the block is angled or straight -- if it is straight, then the depth is consistent across the block, but if it is angled then there are points that are higher or lower, which allows us to tell that it is angled. Importantly, we cannot simply use the depth sensor on the end-effector camera as the camera is too close to where the block would be gripped and thus out of the supported range for the camera. Instead, we use a pre-trained RGB-to-depth neural network model (Depth Anything V2) to produce a depth map which uses the RGB image from the same camera. See Appendix~\ref{appendix:depth_anything} for more details.

\noindent Additionally, after we pick up a block we bring it up to the top camera to determine whether there are any blocks that are attached to it (i.e. the block was improperly detected as not connected, our separation maneuver failed to separate blocks, or the block attached to another block as it was getting picked up). We determine this by simply bringing the block directly under the top camera of the robot and filtering for blocks that appear very large (take up 150 px) that are close to the center of the frame. This allows us to determine how many blocks are close to the top camera, which correspond to the number of blocks the robot has picked up. Note that we decided to use the top-camera for this detection as 1) it allowed us to use our existing, thoroughly-tested block detection algorithm and 2) the end effector camera has a limited and unpredictable view of the gripping area (especially for certain edge cases).

\noindent If the robot detects any mistakes in the grip or connections of the block it has picked up, it is able to recover from it (which will be explained later in this paper).
